\section{Every stratum runs a clock, on its own local group}
\label{sec:strata}

Sections~\ref{sec:clock} and \ref{sec:causal} measure the unit sub-grid. At a field that is
almost the whole table; at $n = 120$ it is $7.1\%$ of it. This section measures the rest, and
it is the paper's central result.

\subsection{What the algebra says to measure}
\label{sec:strata-algebra}

Theorem~\ref{thm:stratum} determines the measurement completely, which is what makes it a
test rather than a search. Fix a stratum $J_{d} \times J_{e}$ and let $m, w, q, G_{m}$ be as in
\eqref{eq:mw} and \eqref{eq:localgroup}. By \eqref{eq:fibration} the target factors through the
fibration, so the measurement is fixed in three steps:
\begin{equation}
  \label{eq:pipeline}
  \underbrace{L\big|_{J_{d} \times J_{e}}}_{\text{block of logits}}
  \;\xrightarrow{\ \iota_{d}^{-1} \times \iota_{e}^{-1}\ }\;
  \underbrace{\widetilde{L}(u, v)}_{\text{unit coordinates}}
  \;\xrightarrow{\ \pi_{d} \times \pi_{e}\ }\;
  \underbrace{\overline{L}(\bar{u}, \bar{v})}_{\text{collapsed onto } G_{m} \times G_{m}}
  \;\xrightarrow{\ \eqref{eq:transform}\ }\;
  \widehat{\overline{L}}_{\kappa, \kappa'} ,
\end{equation}
after which we ask whether the energy sits on $\Delta$ of \eqref{eq:diagonal}. Each arrow is
forced: $\iota$ by Proposition~\ref{prop:torsor}, $\pi$ by \eqref{eq:fibration}, and $\Delta$
by consequence~3 of Theorem~\ref{thm:stratum}.

No key-frequency detector is used anywhere in this section. The ablated set in
\eqref{eq:ablation} is $K = \Delta$, the \emph{whole diagonal}, fixed in advance by
\eqref{eq:diagonal} with $|\Delta| = |G_{m}|$. The obvious objection to any ablation result is
circularity: choosing a set by ablation rank and then asking a permutation test whether that
set is unusually damaging. Here the set cannot be
chosen, because the algebra names it.

The fibre-constancy of the target is asserted at every stratum of all six moduli by the
analysis code rather than assumed.

\subsection{Result}
\label{sec:strata-result}

Criteria were committed before any per-stratum number existed. We measure $32$ strata across
six moduli, $157$ stratum-run measurements in total. They span $12$ distinct local groups
($\mathbb{Z}_{10}$, $\mathbb{Z}_{16}$, $\mathbb{Z}_{20}$, $\mathbb{Z}_{100}$,
$\mathbb{Z}_{110}$, $\mathbb{Z}_{112}$, $\mathbb{Z}_{2} \times \mathbb{Z}_{10}$,
$\mathbb{Z}_{4} \times \mathbb{Z}_{10}$, $\mathbb{Z}_{2}^{2} \times \mathbb{Z}_{4}$,
$\mathbb{Z}_{2} \times \mathbb{Z}_{4} \times \mathbb{Z}_{10}$,
$\mathbb{Z}_{2}^{3} \times \mathbb{Z}_{4}$ and $\mathbb{Z}_{6} \times \mathbb{Z}_{16}$),
and every one carries the same circuit (Figure~\ref{fig:strata}).

\paragraph{Which strata are measurable, and why the cut cannot launder a failure.} A stratum
is \emph{measurable} when two structural conditions fixed in the pre-registration both hold.
Its local group satisfies $|G_{m}| \ge 10$ (S-b: below that the eight-bin family cannot
discriminate), and the split leaves it at least $100$ held-out cells (S-c: fewer cannot
separate a clock from memorisation). Blocks with $d = n$ or $e = n$ are excluded outright,
since $J_{n} = \{0\}$ makes them identically zero. Both conditions are read off the
algebra of $n$ and off the train/test split, never off a model's output, so the cut is
settled before any accuracy is computed and an exclusion can never stand in for a failure.
For G0 in particular, were the cut a function of the outcome, G0 could not fail by
construction.

\begin{center}
\small
\begin{tabular}{l p{0.28\textwidth} l p{0.34\textwidth}}
  \toprule
  & criterion & threshold & observed \\
  \midrule
  G0 & held-out accuracy per block & $\ge 0.95$ & $1.0000 \pm 0.0000$ at every measurable stratum \\
  G1 & permutation $p$, diagonal vs $10{,}000$ sets & $p < 0.01$ in $\ge 4/5$ seeds & $0$ of $10{,}000$ draws in $133/157$; $\hat{p} \le 0.0023$ in all $157$ \\
  G2 & excluded loss / baseline & $\ge 100\times$ & $3.98 \times 10^{2}$ to $2.71 \times 10^{8}$ \\
  G3 & fibre residual \eqref{eq:resid} vs permuted control & $\le 0.5 \times$ control & real $0.000$--$0.132$, control $0.396$--$0.929$ \\
  G4 & diagonal energy share $D$ of \eqref{eq:dshare} & $R \ge 5.0$ & $D = 0.837$--$0.997$ \\
  G5 & tuned neurons $-$ per-stratum shuffle & $\ge 0.20$ & $85$ of $87$ scored; $0.193$--$1.000$ vs shuffle $0.0000$ \\
  G6 & stratum sets differ by the algebra of $n$ & chain vs lattice & see below \\
  \bottomrule
\end{tabular}
\end{center}

G5 is defined only where the tuning statistic is, which is $87$ of the $157$ measurable
stratum-runs (Section~\ref{sec:strata-notclaimed}). Two of those $87$ fall below its
threshold: $n = 165$'s $J_{5} \times J_{15}$ and $J_{15} \times J_{5}$, both at $0.1934$
against a $0.20$ bar. This is the lower end of the range reported above and is not a pass.

Measuring the non-unit strata raises the share of the multiplication table the mechanism
accounts for. Writing $\mathcal{M}$ for the set of measurable strata, that share is
\begin{equation}
  \label{eq:coverage}
  \mathrm{cov}(\mathcal{M}) \;=\;
  \frac{1}{n^{2}} \sum_{(d,e) \in \mathcal{M}} |J_{d}| \, |J_{e}|
  \;=\; \frac{1}{n^{2}} \sum_{(d,e) \in \mathcal{M}} \varphi(n/d)\,\varphi(n/e),
\end{equation}
using $|J_{d}| = \varphi(n/d)$ from \eqref{eq:jclass}. It rises from the unit stratum alone
($\mathcal{M} = \{(1,1)\}$, giving $\varphi(n)^{2}/n^{2}$) to:

\begin{center}
\begin{tabular}{lccl}
  \toprule
  $n$ & unit stratum alone & measured here & local groups \\
  \midrule
  $113$ & $98.2\%$ & $98.2\%$ & $\mathbb{Z}_{112}$ (there is no non-unit stratum) \\
  $121 = 11^{2}$ & $82.6\%$ & $\mathbf{97.7\%}$ & $\mathbb{Z}_{110} \rhd \mathbb{Z}_{10}$ \\
  $125 = 5^{3}$ & $64.0\%$ & $\mathbf{89.6\%}$ & $\mathbb{Z}_{100} \rhd \mathbb{Z}_{20}$ \\
  $119 = 7 \cdot 17$ & $65.1\%$ & $\mathbf{88.6\%}$ & $\mathbb{Z}_{6} \times \mathbb{Z}_{16}$, $\mathbb{Z}_{16}$ \\
  $120$ & $7.1\%$ & $\mathbf{23.1\%}$ & $\mathbb{Z}_{2}^{3} \times \mathbb{Z}_{4}$, $\mathbb{Z}_{2}^{2} \times \mathbb{Z}_{4}$ \\
  $165 = 3 \cdot 5 \cdot 11$ & $23.5\%$ & $\mathbf{82.3\%}$ & $\mathbb{Z}_{2} \times \mathbb{Z}_{4} \times \mathbb{Z}_{10}$, and three more \\
  \bottomrule
\end{tabular}
\end{center}

% STE descriptive
\begin{figure}[t]
  \centering
  \includegraphics[width=\textwidth]{../figures/gate2_strata.pdf}
  \caption[Every measurable stratum, and the criteria it meets.]{\textbf{Every measurable
    stratum, and the criteria it meets.} Each row is one stratum $J_{d} \times J_{e}$, and the
    rows are grouped by modulus. $113$ contributes a single unit stratum, $121$ and $125$ the
    nested $p$-adic chains, $119$ and $165$ divisor lattices, and $120$ the widest lattice.
    Each row carries the stratum's local group $G_{m} = (\mathbb{Z}/(n/m)\mathbb{Z})^{\times}$
    and its score on each criterion of the table above. There are twelve distinct groups
    across the $32$ strata, from $\mathbb{Z}_{10}$ to
    $\mathbb{Z}_{2} \times \mathbb{Z}_{4} \times \mathbb{Z}_{10}$. The point of the figure is
    that the local group changes from row to row while the pattern of met criteria does
    not.\par
    The algebra of $n$ fixes how many strata there are and what group each carries, not
    whether a clock runs on it. Strata excluded by the measurability cut are not drawn, and
    the cut is structural and is defined above. The figure regenerates from
    \texttt{src/viz/plots.py::gate2\_strata}.}
  \label{fig:strata}
\end{figure}
% STE end

G6 is the second half of the claim. The stratum sets differ by the algebra of $n$ in
the way the algebra predicts: $121$ and $125$ give \emph{nested chains}, the $p$-adic
filtration; $119$, $120$ and $165$ give \emph{non-chain divisor lattices}; and $113$ gives a
single stratum. The mechanism does not change. The stratification does.

\subsection{The pre-registered primary criterion does not discriminate}
\label{sec:strata-falsifier}

The pre-registration named its own falsifier: a failed run must not show a local clock. It
fired. It is scored over $183$ grokked and $20$ failed stratum-measurements. These are the
$157$ measurable stratum-runs of the PyTorch sweep above, plus a further $26$ from the
from-scratch engine arm of Section~\ref{sec:strata-engine}, a larger population than the $157$
used for the criteria table.

% STE descriptive
\begin{figure}[t]
  \centering
  \includegraphics[width=\textwidth]{../figures/F9_discrimination.pdf}
  \caption[The pre-registered primary criterion does not separate, and the three that
    do.]{\textbf{The pre-registered primary criterion does not separate, and the three that
    do.} Each row is one Gate 2 criterion, scored on every measurable stratum of the grokked
    runs (filled, $183$) and the failed runs (open, $20$). The count is printed at each mark,
    and the rule between them is the criterion's power to discriminate.
    Rows are ordered by that gap. G0, G2 and G5 separate at $1/20$, $1/20$ and $1/14$. One
    failed measurement passes each, and it is the same one, named in the text.\par
    G3 and G4 separate partly. G1, the permutation $p$ named primary in the pre-registration,
    passes on $20/20$ failed measurements. Its two marks coincide, so the gap is zero. The
    verdict rests on G0 $\wedge$ G2 $\wedge$ G5. Counts differ by row, because G3 and G5 are
    undefined on strata with no fibre to collapse and on runs without saved activations. A
    row is scored only where it is defined.\par
    The figure regenerates from \texttt{src/viz/plots.py::gate2\_discrimination}, scored by
    \texttt{scripts/gate2\_discriminate.py}.}
  \label{fig:discrimination}
\end{figure}
% STE end

Once looked at, the explanation is not subtle: $30\%$ of every block is training data, so a model
that has merely memorised its training cells still has logits that are a function of
$u \cdot v$ there. G1 tests whether the logits are a function of the product, and
memorisation satisfies that. It is a necessary condition, not evidence. It was named primary
because it is the statistic this project had used elsewhere. That is the kind of inherited
authority that should have been checked against a failed-run control before it was promoted,
and now is.

The verdict stands, and its weight rests on G0 $\wedge$ G2 $\wedge$ G5, each of which
passes on exactly one failed measurement and no more. We do not write that all six criteria
passed.

That one measurement is worth naming rather than rounding away. It is the unit
stratum of $n = 49$, seed $2$, of the extended sweep, and it is the only measurable stratum
that run has. Its window-median test accuracy is $0.8908$, which clears our FAILED threshold
of $0.90$ by $0.009$, and its final ten logged samples rise, with one small dip, from $0.842$ to
$0.963$. It is a run still learning when the step budget ended, not one that plateaued having
learned nothing. On that stratum it reaches $0.959$ held-out accuracy and an
excluded/baseline ratio of $302$, so it passes G0, G2, G4 and G5 exactly as a grokked
measurement would. We place it in the failed arm because that is where our three-state
endpoint rule puts it. A FAILED threshold $0.01$ \emph{lower}, at $0.89$, would reclassify it as a near-grok, remove it from the arm, and restore the
$0/19$ this paper previously reported. The threshold is not tuned to this measurement: the
$0.90$ floor is inherited from the excursion analysis that retired our own de-grokking claim
(Appendix~\ref{app:retractions}), fixed a week before this arm was scored, and we keep it. The claim is therefore that G0, G2 and
G5 separate the two arms on every measurement but one borderline run, not that they separate
completely.

This is a narrowing, not a rescue, and the direction is the whole argument. The
decision rule was fixed in the pre-registration before the data existed: the fourth branch is
declared if G1 $\wedge$ G2 hold at every measurable non-unit stratum at two or more moduli and
G6 shows the stratum sets differ. That rule was met on its own terms: G1 \emph{held},
at $99\%$ of grokked stratum-measurements. What fired is a \emph{separate} clause of the same
document, the falsifier requiring that a failed run not show a local clock, and it is a finding
about what G1 can discriminate rather than about whether the rule was satisfied. The move we
then make, resting the weight on G0 $\wedge$ G2 $\wedge$ G5, \emph{discards} the
criterion our own pre-registration named primary and leaves a strictly smaller evidence base
than the one we registered. A post-hoc rescue selects criteria to save a claim; this
deselects one, after it passed, because it turned out not to separate. The same direction runs
through Appendix~\ref{app:retractions}: all five withdrawn claims were positive findings of
our own, three of them died to a size confound we went looking for, and not one retraction
replaced a claim with a stronger one.

\subsection{The same clock, or one clock each}
\label{sec:strata-sameness}

Every criterion above is computed inside one stratum. That cannot separate \emph{one
mechanism, indexed per stratum} from \emph{a private sub-circuit per stratum}: both put a
diagonal clock in every block. The discriminating question is whether two strata that share
a local group select the \emph{same characters} of it. The instrument for it is already in
this paper: \eqref{eq:keyagree} compares two independently recovered key sets element
for element on the unit stratum. Here it is applied between strata, pre-registered before
any cross-stratum number existed.

Which pairs can carry the test is fixed by the algebra, before any model is opened.
Two strata are comparable when their local groups are identical. $\mathcal{J}$-blocks
$(d,e)$ and $(e,d)$ are transposes of one \emph{commutative} task, so their diagonal spectra
agree for a reason that has nothing to do with a shared circuit; they are reported
separately and are never support. What remains is $28$ pairs at $n = 165$, $2$ at $n = 120$
and $2$ at $n = 119$, and \emph{none} at $n = 121$ or $n = 125$, whose only same-group
pair is a transpose. $n = 165$'s $G_{11}$, eight strata on one local group, carries the test.
Across five seeds that is $160$ pairs.

\begin{center}
\begin{tabular}{llcc}
  \toprule
  & statistic & observed & pre-registered \\
  \midrule
  \textbf{agreement} & key sets \emph{exactly} equal & $\mathbf{158/160 = 0.9875}$ & $\ge 0.80$ \\
  \textbf{null} & Monte-Carlo, key sets redrawn at the observed sizes & $\mathbf{1.0 \times 10^{-4}}$ & $< 0.01$ \\
  \textbf{control} & cell-shuffle, sizes kept, structure destroyed & $\mathbf{0/160}$ & $< 0.30$ \\
  transpose & descriptive only, never support & $54/54$ & --- \\
  \bottomrule
\end{tabular}
\end{center}

\noindent
Mean Jaccard overlap is $0.9982$ and the median key set has $2$ characters. The two
disagreements are both in $n = 165$ seed $1$, which reads $26/28$; the other four seeds read
$28/28$. The null is the floor a $10{,}000$-draw design can express, and it is not a formality:
at $G_{11}$ the folded spectrum has $5$ bins, so two independent $2$-character sets coincide
by chance with probability $1/\binom{5}{2} = 0.10$. The shuffle control does not fail by
selecting \emph{different} characters: $0$ of $160$ shuffled blocks has any key set at all,
so the shuffle removes the clock rather than relabelling it.

What this does not have is a failed-run control. The design called for the same statistic on
models that never generalised. There are exactly \emph{two} scorable pairs of that kind in the
whole project, and both come from one run at test accuracy $0.7502$. That model learned most
of the task, and our three-state rule nonetheless buckets it with the $0.146$ memorisers. Both agree, so the control
reads $2/2$ against its own $0.50$ bar; its own null reads $p = 1.0$, because two pairs
cannot discriminate anything. Every genuinely failed run either has no comparable pair by
algebra ($n = 49$, $63$, $125$) or produces \emph{no key set at all} ($n = 54$ seed $1$, at
accuracy $0.2733$). The control is underpowered rather than passing, and this section's claim
rests on the null and the shuffle control instead. The same shape, a statistic that is load-bearing elsewhere passing on failures here, is what
Section~\ref{sec:strata-falsifier} reports for the pre-registered primary criterion. It is
why we score every criterion on a failure when a failure exists.

\subsection{The single-seed precursor}
\label{sec:strata-c7}

The torsor coordinate was first sighted at one seed, before the pre-registration existed, by
indexing a non-regular $\mathcal{J}$-class not by a group law it lacks but by the unit action,
and taking the product-character transform in that coordinate. Every class beat its
random-basis control, regular or not: $J_{11}$ at $n = 121$ scored $0.345$ against $0.124$
($2.79\times$), $J_{5}$ at $n = 125$ scored $0.526$ against $0.156$ ($3.37\times$), and
$J_{2}$ at $n = 120$ scored $0.476$ against $0.162$ ($2.93\times$) (1 seed). That
observation is the origin of this section, and the section supersedes it: the same statement,
measured at five seeds across $26$ non-unit strata under criteria committed in advance.

An aggregate drawn from that single-seed table, that regular classes carry more structure
than non-regular ones, was retracted. It was confounded by class size, $J_{1}$ being always
regular and always largest. Size-matched within one model at $n = 120$, the only place such
pairs exist, the two overlap. Appendix~\ref{app:retractions} carries it with the other
retractions.

\subsection{Independent replication}
\label{sec:strata-engine}

The same verdict holds on our from-scratch engine at three moduli ($119$, $121$, $125$), two
to three seeds each: every measurable non-unit stratum carries a local clock. The two
implementations disagree about sparsity \emph{magnitude} for the reason given in
Section~\ref{sec:impl}, and agree about the stratified clock. One miss is reported:
the $n = 119$ unit stratum fails G1 in one of two seeds.

\subsection{What is not claimed}
\label{sec:strata-notclaimed}

\paragraph{G1, G2 and G4 are partly entailed by correctness.} By \eqref{eq:fibration} any
\emph{correct} function on a stratum is a function of $\mathrm{dlog}(\bar{u}) +
\mathrm{dlog}(\bar{v})$, hence supported on $\Delta$. So a perfectly confident lookup table
already reads $D = 1.000$, $\hat{p}$ at the floor and fibre residual $0$, and the analysis code
plants exactly that and measures it. What these criteria establish is that the logit tensor
depends on $(u,v)$ essentially only through $u \cdot v \bmod n/m$, i.e.\ that the computation
has collapsed onto the local ring. A lookup-table fit of $r^{2} = 0.012$--$0.074$ says the
logits are far from a scaled one-hot, so the collapse is a real statement about the tensor.
But a smooth function of $u \cdot v$ would also read $D \approx 1$, and these criteria cannot
by themselves say the mechanism is a \emph{character} circuit.

Sameness across strata is measured in Section~\ref{sec:strata-sameness}, and the limit of
that measurement is the absence of a failed-run control, not the absence of the
measurement. Every other statistic in this section is computed \emph{inside} one stratum,
against that stratum's own local group. A model that routed each stratum to a private
sub-circuit, each correct on its own block, would produce all of them.

G5 is the criterion that says ``clock''. Single-frequency MLP neurons in local-group
coordinates are entailed by nothing, and the shuffle control reads $0.0000$ at every stratum.
The tuning statistic as pre-registered is cyclic-only, so G5 covers $121$, $125$, $119$'s
sub-strata and $165$'s $\mathbb{Z}_{10}$ strata. The non-cyclic strata are covered by a
product-character version reading $0.182$--$0.877$ against a shuffle control of $0.0000$,
which was added after the pre-registration, is exploratory, and was never scored as G5.

The model does not solve every stratum. On held-out cells, $15$ stratum-runs fall to
$0.379$--$0.833$. Every one is below the measurability cut, and none affects the verdict. An
earlier internal record of ours claimed every stratum at accuracy $1.0000$, but that was the
full grid including training cells, and it is corrected here. The degradation is also not
evidence about the algebra. Every failing stratum has $|G_{m}| \le 8$ \emph{and} at most $64$
cells, and those are structurally coupled. So it cannot separate ``the local group is too
small to carry a clock'' from ``too few training examples''.

$n = 113$ contributes nothing to the stratum claim. A field has one stratum. It is
the positive control and that is all.
